\documentclass[11pt]{article}
\usepackage{mathblatt}
\newcommand{\zeichenplatz}[1][5]{\begin{tikzpicture}\path (0,0) rectangle (#1,#1);\fill (#1/2,#1/2) circle (0.06) node[below right]{$M$};\end{tikzpicture}}
\begin{document}
\blattkopf{Kreis}{Kennst du schon}
\weit
\einheitenkopf[zone][Kennst du schon]{Das kennst du schon}
\begin{aufgabe}{Ich kann einen Kreis mit dem Zirkel zeichnen.}
Zeichne um $M$ den Kreis mit dem Zirkel.
\begin{teile}
\teil Radius $r = 2\,\text{cm}$ \quad \teil Durchmesser $d=3$ cm
\end{teile}

\zeichenplatz\hfill\zeichenplatz\hfill\zeichenplatz

\end{aufgabe}
\begin{aufgabe}{Probe Kreissektor.}
\kreissektor{90}{}\hspace{5mm}\kreissektor{235}{}

\begin{kreis}[1.3]\mittelpunkt{M}\radius{40}{1}\durchmesser{160}{2}\sehne{200}{280}{3}\end{kreis}
\hspace{1cm}
\begin{kreis}[1.2]\mittelpunkt{}\sektor{0}{90}{}\radius{0}{3\,\text{cm}}\end{kreis}
\hspace{1cm}
\begin{kreis}[1.2]\mittelpunkt{}\sektor{0}{150}{150^\circ}\radius{0}{7\,\text{m}}\end{kreis}

\winkel[4]{65}{}

\winkelstrahl[6]
\end{aufgabe}
\begin{aufgabe}{Tabelle.}
\sachtabelle{lcccc}{ & Radius $r$ & Durchmesser $d$ & Umfang $u$ & Fläche $A$}{a) & $6\,$m & \leerzelle & \leerzelle & \leerzelle \\ b) & \leerzelle & $9$ cm & \leerzelle & \leerzelle}
\begin{teile}\teil Strecke 1: \kreuz{Radius}\kreuz{Durchmesser}\kreuz{keins von beiden}
\teil $u=4\cdot a$ \feld{u} \leerfeld[\%]
\end{teile}
\kreisleer[2]
\end{aufgabe}
\end{document}
